\setcounter{section}{3}


  \zwischenueberschrift{Die Zariski-Topologie}


\bild{ \begin{center}


\includegraphics[width=5.5cm]{ \bildeinlesungjpg {Oscar_Zariski} {jpg} }


\end{center}


\bildtext {
Oscar Zariski (1899-1986)
} }


\bildlizenz { Oscar Zariski.jpg } {} {} {Commons} {} {}


In

Proposition 2.8

haben wir gezeigt, dass die affin-algebraischen Teilmengen eines affinen Raumes die Axiome für abgeschlossene Mengen einer
\definitionsverweis {Topologie}{}{}
erfüllen. Diese Topologie nennt man die
Zariski
-Topologie.


\inputdefinition


{}


{


In einem
\definitionsverweis {affinen Raum}{}{}


\mathl{{ {\mathbb A}_{  K }^{  n   } }}{} versteht man unter der \definitionswort {Zariski-Topologie}{} diejenige
\definitionsverweis {Topologie}{}{,}
bei der die
\definitionsverweis {affin-algebraischen Mengen}{}{}
als abgeschlossen erklärt werden.


}


Die offenen Mengen der Zariski-Topologie sind also die Komplemente der affin-algebraischen Mengen. Sie werden zu einem Ideal ${\mathfrak a}$ mit


\mavergleichskettedisp


{\vergleichskette


{ D({\mathfrak a})
}


{ =} { { {\mathbb A}_{ K }^{ n  } }  \setminus V({\mathfrak a})
}


{ } {
}


{ } {
}


{ } {
}


}
{}{}{}
bezeichnet. Die Zariski-Topologie weicht sehr stark von anderen Topologien ab, insbesondere von solchen, die durch eine Metrik gegeben sind. Insbesondere ist die Zariski-Topologie nicht
\definitionsverweis {hausdorffsch}{}{.}
Generell kann man sagen, dass die offenen Mengen
\zusatzklammer {außer der leeren Menge} {} {}
in der Zariski-Topologie sehr groß sind
\zusatzklammer {siehe

Aufgabe 3.20
} {} {,}
während die abgeschlossenen
\zusatzklammer {also die affin-algebraischen Mengen} {} {}
sehr dünn sind
\zusatzklammer {außer dem ganzen Raum selbst} {} {.}


\inputbeispiel{}


{


Die Zariski-Topologie auf der affinen Geraden


\mathl{{\mathbb A}^{1}_{ K }}{} über einem
\definitionsverweis {Körper}{}{}
$K$ lässt sich einfach beschreiben. Als
\zusatzklammer {Zariski} {} {-}abgeschlossene Teilmenge haben wir zunächst einmal die gesamte affine Gerade, die durch


\mathl{V(0)}{} beschrieben wird. Alle anderen abgeschlossenen Teilmengen werden durch


\mathl{V({\mathfrak a})}{} mit


\mavergleichskette


{\vergleichskette


{ {\mathfrak a}
}


{ \neq }{ 0
}


{  }{
}


{    }{
}


{   }{
}


}
{}{}{}
beschrieben. Da


\mathl{K[X]}{} ein
\definitionsverweis {Hauptidealbereich}{}{}
ist, kann man sogar


\mathbed {{\mathfrak a} =(f)} {}


 {f \neq 0} {}


 {} {} {} {,}
ansetzen. Die zugehörige Nullstellenmenge besteht also aus endlich vielen Punkten. Andererseits ist jeder einzelne Punkt $P$ mit der Koordinate $a$ die einzige Nullstelle des linearen Polynoms


\mathl{X-a}{,} also ist


\mavergleichskette


{\vergleichskette


{ \{P\}
}


{ = }{ V(X-a)
}


{  }{
}


{    }{
}


{   }{
}


}
{}{}{}
Zariski-abgeschlossen. Eine endliche Ansammlung von Punkten


\mathl{P_1 , \ldots ,  P_k}{} mit den Koordinaten


\mathl{a_1 , \ldots , a_k}{} ist die Nullstellenmenge des Polynoms


\mathl{(X-a_1) \cdots (X-a_k)}{.} Die Zariski-abgeschlossenen Mengen der affinen Geraden bestehen also aus allen endlichen Teilmengen
\zusatzklammer {einschließlich der leeren} {} {}
und der gesamten Menge.


}


\bild{ \begin{center}


\includegraphics[width=5.5cm]{ \bildeinlesungjpg {Lineline} {jpg} }


\end{center}


\bildtext {} }


\bildlizenz { Lineline.jpg } {} {Astur1} {Commons} {PD} {}


\inputbeispiel{}


{


Jeder Punkt


\mavergleichskettedisp


{\vergleichskette


{P
}


{ =} { (a_1 , \ldots , a_n)
}


{ \in} { { {\mathbb A}_{  K }^{  n   } }
}


{ } {
}


{ } {
}


}
{}{}{}
ist
\definitionsverweis {Zariski-abgeschlossen}{}{,}
und zwar ist


\mavergleichskettedisp


{\vergleichskette


{P
}


{ =} { V(X_1-a_1,X_2-a_2 , \ldots , X_n- a_n )
}


{ } {
}


{ } {
}


{ } {
}


}
{}{}{.}
Punkte sind
\zusatzklammer {neben der leeren Menge und dem gesamten Raum} {} {}
die einfachsten affin-algebraischen Mengen. Das Ideal
\zusatzklammer {genannt \stichwort {Punktideal} {}} {} {}


\mathl{(X_1-a_1, X_2-a_2 , \ldots , X_n- a_n )}{} ist
\definitionsverweis {maximal}{}{,}
siehe

Aufgabe 2.12
.


}


\bild{ \begin{center}


\includegraphics[width=5.5cm]{ \bildeinlesungpng {IntersectingPlanes} {png} }


\end{center}


\bildtext {Der Schnitt von zwei und} }


\bildlizenz { IntersectingPlanes.png } {} {Stib} {en.wikipedia.org} {CC-BY-SA-3.0} {}


\bild{ \begin{center}


\includegraphics[width=5.5cm]{ \bildeinlesungpng {Secretsharing-3-point} {png} }


\end{center}


\bildtext {von drei Ebenen} }


\bildlizenz { Secretsharing-3-point.png } {} {Stib} {en.wikipedia.org} {CC-BY-SA-3.0} {}


Nach

Proposition 2.8

 (3)

sind somit alle endlichen Teilmengen des affinen Raumes Zariski-abgeschlossen und somit sind die Komplemente


\mathl{{ {\mathbb A}_{ K }^{ n  } } \setminus E}{,} wenn $E$ eine endliche Punktemenge bezeichnet, Zariski-offen. Ferner ist zu einer rationalen Funktion


\mathl{P/Q}{} mit


\mavergleichskette


{\vergleichskette


{ P,Q
}


{ \in }{ K[X_1 , \ldots , X_n]
}


{  }{
}


{    }{
}


{   }{
}


}
{}{}{}
und


\mavergleichskette


{\vergleichskette


{ Q
}


{ \neq }{ 0
}


{  }{
}


{    }{
}


{   }{
}


}
{}{}{}
der Definitionsbereich, nämlich $D(Q)$, offen.


  \zwischenueberschrift{Das Verschwindungsideal}


\inputdefinition


{}


{


Es sei


\mavergleichskette


{\vergleichskette


{T
}


{ \subseteq }{ { {\mathbb A}_{  K }^{  n   } }
}


{  }{
}


{    }{
}


{   }{
}


}
{}{}{}
eine Teilmenge. Dann nennt man


\mathdisp {\left\{  F \in  K[X_1 , \ldots ,  X_n]   \mid   F(P) = 0 \text { für alle } P \in T        \right\}} {  }


das \definitionswort {Verschwindungsideal}{} zu $T$. Es wird mit


\mathl{\operatorname{Id} \,(T)}{} bezeichnet.


}


Es handelt sich dabei in der Tat um ein Ideal: Wenn


\mavergleichskette


{\vergleichskette


{ F(P)
}


{ = }{ 0
}


{  }{
}


{    }{
}


{   }{
}


}
{}{}{}
und


\mavergleichskette


{\vergleichskette


{ G(P)
}


{ = }{ 0
}


{  }{
}


{    }{
}


{   }{
}


}
{}{}{}
für alle


\mavergleichskette


{\vergleichskette


{ P
}


{ \in }{ T
}


{  }{
}


{    }{
}


{   }{
}


}
{}{}{}
ist, so gilt dies auch für die Summe


\mathl{F+G}{} und für jedes Vielfache $HF$.


Damit haben wir zwei Zuordnungen in entgegengesetzte Richtung: Einer Teilmenge im affinen Raum wird das Verschwindungsideal zugeordnet und einem Ideal im Polynomring das zugehörige Nullstellengebilde. Wir interessieren uns dafür, inwiefern sich Ideale und Nullstellengebilde entsprechen.


\inputbeispiel{}


{


Das
\definitionsverweis {Verschwindungsideal}{}{}
zur leeren Menge ist das
\definitionsverweis {Einheitsideal}{}{,}
da es keinen Punkt gibt, auf dem die Nullstellenbedingung überprüft werden müsste.


Das Verschwindungsideal zum Gesamtraum


\mathl{{ {\mathbb A}_{  K }^{  n   } }}{} hängt vom Körper ab. Wenn dieser unendlich ist, so gibt es nur das Nullpolynom, das überall verschwindet, und folglich ist das Verschwindungsideal gleich dem Nullideal. Dies folgt aus

Aufgabe 3.18
.


Ist hingegen der Körper endlich mit $q$ Elementen, so ist


\mavergleichskette


{\vergleichskette


{ x^q -x
}


{ = }{ 0
}


{  }{
}


{    }{
}


{   }{
}


}
{}{}{}
für jedes


\mavergleichskette


{\vergleichskette


{ x
}


{ \in }{ K
}


{  }{
}


{    }{
}


{   }{
}


}
{}{}{.}
Also verschwindet das Polynom


\mathl{X^q-X}{} auf jedem Punkt der affinen Geraden und gehört somit zum Verschwindungsideal der affinen Geraden. In höherer Dimension ist das Verschwindungsideal von


\mathl{{ {\mathbb A}_{  K }^{  n   } }}{}   gleich


\mathl{(X_1^q-X_1,X_2^q-X_2 , \ldots , X_n^q-X_n)}{.}


}


\inputbeispiel{}


{


Es sei


\mavergleichskette


{\vergleichskette


{ P
}


{ = }{ (a_1 , \ldots , a_n)
}


{ \in }{ { {\mathbb A}_{  K }^{  n   } }
}


{    }{
}


{   }{
}


}
{}{}{.}
Dann ist das
\definitionsverweis {Verschwindungsideal}{}{}


\mathl{\operatorname{Id}  \, (P)}{} gleich dem Ideal


\mathl{(X_1-a_1 , \ldots ,  X_n -a_n)}{.} Zunächst ist klar, dass die linearen Polynome


\mathl{X_i-a_i}{} im Punkt $P$ verschwinden
\zusatzklammer {wegen


\mavergleichskettek


{\vergleichskettek


{ (X_i-a_i)(P)
}


{ = }{ a_i -a_i
}


{ = }{ 0
}


{    }{
}


{   }{
}


}
{}{}{}} {} {.}
Damit gehört auch das von diesen Polynomen erzeugte Ideal zum Verschwindungsideal. Es sei umgekehrt $F$ ein Polynom mit


\mavergleichskette


{\vergleichskette


{ F(P)
}


{ = }{ 0
}


{  }{
}


{    }{
}


{   }{
}


}
{}{}{.}
Wir schreiben $F$ in den \anfuehrung{neuen Variablen}{}


\mathdisp {\tilde{X}_1 = X_1-a_1 , \ldots ,  \tilde{X}_n = X_n -a_n} { , }


indem wir $X_i$ durch


\mathl{X_i-a_i+a_i}{} ersetzen. In den neuen Variablen sei


\mavergleichskette


{\vergleichskette


{ F
}


{ = }{ \sum_\nu b_\nu \tilde{X}^\nu
}


{  }{
}


{    }{
}


{   }{
}


}
{}{}{.}
Dieses Polynom besteht aus der Konstanten $b_0$, in jedem anderen Monom kommt mindestens eine Variable vor. Also können wir


\mavergleichskettedisp


{\vergleichskette


{F
}


{ =} { F_1 \tilde{X}_1 + \cdots + F_n \tilde{X}_n + c
}


{ } {
}


{ } {
}


{ } {
}


}
{}{}{}
mit gewissen Polynomen $F_i$ schreiben. Daher ist


\mavergleichskette


{\vergleichskette


{ F(P)
}


{ = }{ c
}


{ = }{ 0
}


{    }{
}


{   }{
}


}
{}{}{}
und


\mavergleichskette


{\vergleichskette


{ F
}


{ \in }{ (\tilde{X}_1 , \ldots ,  \tilde{X}_n)
}


{  }{
}


{    }{
}


{   }{
}


}
{}{}{.}


}


\inputfaktbeweis


{Affine Varietäten/Verschwindungsideal/Antimonotonie/Fakt}


{Lemma}


{}


{


\faktsituation {Es seien


\mavergleichskette


{\vergleichskette


{V
}


{ \subseteq }{ W
}


{ \subseteq }{ { {\mathbb A}_{  K }^{  n   } }
}


{    }{
}


{   }{
}


}
{}{}{}
Teilmengen.}


\faktfolgerung {Dann gilt für die zugehörigen
\definitionsverweis {
Verschwindungsideale
}{}{}
die Inklusion


\mavergleichskettedisp


{\vergleichskette


{ \operatorname{Id}  \, (W)
}


{ \subseteq} { \operatorname{Id} \, (V)
}


{ } {
}


{ } {
}


{ } {
}


}
{}{}{.}}


\faktzusatz {}


\faktzusatz {}


}


{


Es sei


\mavergleichskette


{\vergleichskette


{ F
}


{ \in }{ \operatorname{Id} \,(W)
}


{  }{
}


{    }{
}


{   }{
}


}
{}{}{,}
d.h. es ist


\mavergleichskette


{\vergleichskette


{ F(P)
}


{ = }{ 0
}


{  }{
}


{    }{
}


{   }{
}


}
{}{}{}
für alle


\mavergleichskette


{\vergleichskette


{ P
}


{ \in }{ W
}


{  }{
}


{    }{
}


{   }{
}


}
{}{}{.}
Dann ist erst recht


\mavergleichskette


{\vergleichskette


{ F(P)
}


{ = }{ 0
}


{  }{
}


{    }{
}


{   }{
}


}
{}{}{}
für alle


\mavergleichskette


{\vergleichskette


{ P
}


{ \in }{ V
}


{  }{
}


{    }{
}


{   }{
}


}
{}{}{.}
Also ist auch


\mavergleichskette


{\vergleichskette


{ F
}


{ \in }{ \operatorname{Id} \,(V)
}


{  }{
}


{    }{
}


{   }{
}


}
{}{}{.}


}


\inputfaktbeweis


{Affine Varietäten/Verschwindungsideal zu Teilmenge/Nullstellengebilde/Beziehung/Fakt}


{Lemma}


{}


{


\faktsituation {}


\faktvoraussetzung {Es sei


\mavergleichskette


{\vergleichskette


{ I
}


{ \subseteq }{ K[X_1 , \ldots ,  X_n]
}


{  }{
}


{    }{
}


{   }{
}


}
{}{}{}
ein
\definitionsverweis {Ideal}{}{}
und sei


\mavergleichskette


{\vergleichskette


{T
}


{ \subseteq }{ { {\mathbb A}_{  K }^{  n   } }
}


{  }{
}


{    }{
}


{   }{
}


}
{}{}{}
eine Teilmenge.}


\faktuebergang {Dann gelten folgende Aussagen.}


\faktfolgerung {\aufzaehlungvier{Es ist


\mavergleichskette


{\vergleichskette


{T
}


{ \subseteq }{ V ( \operatorname{Id}(T))
}


{  }{
}


{    }{
}


{   }{
}


}
{}{}{.}
}{Es ist


\mavergleichskette


{\vergleichskette


{ I
}


{ \subseteq }{ \operatorname{Id} (V(I))
}


{  }{
}


{    }{
}


{   }{
}


}
{}{}{.}
}{Es ist


\mavergleichskette


{\vergleichskette


{ V(I)
}


{ = }{ V ( \operatorname{Id}(V(I)))
}


{  }{
}


{    }{
}


{   }{
}


}
{}{}{.}
}{Es ist


\mavergleichskette


{\vergleichskette


{ \operatorname{Id} (T)
}


{ = }{ \operatorname{Id} (V( \operatorname{Id} (T)  ))
}


{  }{
}


{    }{
}


{   }{
}


}
{}{}{.}
}}


\faktzusatz {}


\faktzusatz {}


}


{


(1). Es sei


\mavergleichskette


{\vergleichskette


{ P
}


{ \in }{ T
}


{  }{
}


{    }{
}


{   }{
}


}
{}{}{}
ein Punkt. Dann verschwindet nach Definition jedes Polynom


\mavergleichskette


{\vergleichskette


{ F
}


{ \in }{ \operatorname{Id}(T)
}


{  }{
}


{    }{
}


{   }{
}


}
{}{}{}
auf $T$, also


\mavergleichskette


{\vergleichskette


{ P
}


{ \in }{ V(  \operatorname{Id} (T))
}


{  }{
}


{    }{
}


{   }{
}


}
{}{}{.}


(2). Es sei


\mavergleichskette


{\vergleichskette


{ F
}


{ \in }{ I
}


{  }{
}


{    }{
}


{   }{
}


}
{}{}{.}
Dann verschwindet $F$ auf ganz $V(I)$ und daher ist


\mavergleichskette


{\vergleichskette


{ F
}


{ \in }{ \operatorname{Id} \, (V(I))
}


{  }{
}


{    }{
}


{   }{
}


}
{}{}{.}


(3). Nach (1), angewandt auf


\mavergleichskette


{\vergleichskette


{ T
}


{ = }{ V(I)
}


{  }{
}


{    }{
}


{   }{
}


}
{}{}{,}
haben wir die Inklusion \anfuehrung{ $\subseteq$ }{.} Nach (2) ist


\mavergleichskette


{\vergleichskette


{ I
}


{ \subseteq }{ \operatorname{Id} \, (V(I))
}


{  }{
}


{    }{
}


{   }{
}


}
{}{}{.}
Wendet man darauf $V(-)$ an, so ergibt sich nach

Lemma 2.7

die andere Inklusion.


(4). Wie (3).


}


\inputbeispiel{}


{


Die Inklusionen in

Lemma 3.8

(1), (2) sind echt. Es sei zum Beispiel


\mavergleichskette


{\vergleichskette


{ T
}


{ \subset }{ {\mathbb A}^{1}_{K}
}


{  }{
}


{    }{
}


{   }{
}


}
{}{}{}
eine unendliche echte Teilmenge
\zusatzklammer {was voraussetzt, dass $K$ unendlich ist} {} {.}
Dann ist


\mavergleichskette


{\vergleichskette


{ \operatorname{Id} (T)
}


{ = }{ 0
}


{  }{
}


{    }{
}


{   }{
}


}
{}{}{,}
und also ist


\mavergleichskette


{\vergleichskette


{ V(0)
}


{ = }{ {\mathbb A}^{1}_{K}
}


{  }{
}


{    }{
}


{   }{
}


}
{}{}{}
echt größer als $T$.


Zu (2). Es sei


\mavergleichskette


{\vergleichskette


{ I
}


{ = }{ (X^2)
}


{  }{
}


{    }{
}


{   }{
}


}
{}{}{,}


\mavergleichskette


{\vergleichskette


{ R
}


{ = }{ K[X]
}


{  }{
}


{    }{
}


{   }{
}


}
{}{}{.}
Dann ist


\mavergleichskette


{\vergleichskette


{ V(I)
}


{ = }{ \{0\}
}


{  }{
}


{    }{
}


{   }{
}


}
{}{}{}
und


\mavergleichskette


{\vergleichskette


{ \operatorname{Id} \,(\{0\})
}


{ = }{ (X)
}


{  }{
}


{    }{
}


{   }{
}


}
{}{}{,}
aber


\mavergleichskette


{\vergleichskette


{ X
}


{ \notin }{ (X^2)
}


{  }{
}


{    }{
}


{   }{
}


}
{}{}{.}
Ein extremeres Beispiel für


\mavergleichskette


{\vergleichskette


{ R
}


{ = }{ \R [X,Y]
}


{  }{
}


{    }{
}


{   }{
}


}
{}{}{}
ist


\mavergleichskette


{\vergleichskette


{ I
}


{ = }{ (X^2+Y^2)
}


{  }{
}


{    }{
}


{   }{
}


}
{}{}{}
mit


\mavergleichskette


{\vergleichskette


{ V(I)
}


{ = }{ \{(0,0)\}
}


{  }{
}


{    }{
}


{   }{
}


}
{}{}{.}
Das Verschwindungsideal zu diesem Punkt ist aber das Ideal


\mathl{(X,Y)}{.}


}


\inputfaktbeweis


{Affiner Raum/Zariski-Topologie/Zariski-Abschluss ist V zu Verschwindungsideal/Fakt}


{Lemma}


{}


{


\faktsituation {Es sei


\mavergleichskette


{\vergleichskette


{T
}


{ \subseteq }{ { {\mathbb A}_{  K }^{  n   } }
}


{  }{
}


{    }{
}


{   }{
}


}
{}{}{}
eine Teilmenge.}


\faktfolgerung {Dann ist der Zariski-Abschluss von $T$ gleich


\mavergleichskettedisp


{\vergleichskette


{ \overline{T}
}


{ =} { V( \operatorname{Id} \,(T)  )
}


{ } {
}


{ } {
}


{ } {
}


}
{}{}{.}}


\faktzusatz {}


\faktzusatz {}


}


{


Die Inklusion


\mavergleichskette


{\vergleichskette


{ T
}


{ \subseteq }{ V( \operatorname{Id} \,(T) )
}


{  }{
}


{    }{
}


{   }{
}


}
{}{}{}
wurde in

Lemma 3.8

 (1)

gezeigt. Da


\mathl{V( \operatorname{Id} \,(T))}{} nach Definition abgeschlossen ist, folgt daraus


\mavergleichskette


{\vergleichskette


{ \overline{T}
}


{ \subseteq }{ V( \operatorname{Id} \,(T))
}


{  }{
}


{    }{
}


{   }{
}


}
{}{}{.}


Es sei umgekehrt


\mavergleichskette


{\vergleichskette


{ P
}


{ \in }{ V( \operatorname{Id} \,(T))
}


{  }{
}


{    }{
}


{   }{
}


}
{}{}{}
und sei


\mavergleichskette


{\vergleichskette


{ P
}


{ \notin }{ \overline{T}
}


{  }{
}


{    }{
}


{   }{
}


}
{}{}{}
angenommen. Dies bedeutet, dass es eine Zariski-offene Menge $U$ gibt mit


\mavergleichskette


{\vergleichskette


{ P
}


{ \in }{ U
}


{  }{
}


{    }{
}


{   }{
}


}
{}{}{}
und


\mavergleichskette


{\vergleichskette


{ U \cap T
}


{ = }{ \emptyset
}


{  }{
}


{    }{
}


{   }{
}


}
{}{}{.}
Es sei


\mavergleichskette


{\vergleichskette


{ U
}


{ = }{ D( {\mathfrak a} )
}


{  }{
}


{    }{
}


{   }{
}


}
{}{}{.}
Die Bedingung


\mavergleichskette


{\vergleichskette


{ P
}


{ \in }{ U
}


{  }{
}


{    }{
}


{   }{
}


}
{}{}{}
bedeutet, dass es ein


\mavergleichskette


{\vergleichskette


{ G
}


{ \in }{ {\mathfrak a}
}


{  }{
}


{    }{
}


{   }{
}


}
{}{}{}
mit


\mavergleichskette


{\vergleichskette


{ G(P)
}


{ \neq }{ 0
}


{  }{
}


{    }{
}


{   }{
}


}
{}{}{}
geben muss. Es ist dann


\mavergleichskette


{\vergleichskette


{ P
}


{ \in }{ D(G)
}


{ \subseteq }{ U
}


{    }{
}


{   }{
}


}
{}{}{}
und damit


\mavergleichskette


{\vergleichskette


{ T \cap D(G)
}


{ = }{ \emptyset
}


{  }{
}


{    }{
}


{   }{
}


}
{}{}{.}
Also ist


\mavergleichskette


{\vergleichskette


{ T
}


{ \subseteq }{ V(G)
}


{  }{
}


{    }{
}


{   }{
}


}
{}{}{}
und somit


\mavergleichskette


{\vergleichskette


{ G
}


{ \in }{ \operatorname{Id} \,(T)
}


{  }{
}


{    }{
}


{   }{
}


}
{}{}{.}
Wegen


\mavergleichskette


{\vergleichskette


{ G(P)
}


{ \neq }{ 0
}


{  }{
}


{    }{
}


{   }{
}


}
{}{}{}
ergibt sich ein Widerspruch zu


\mavergleichskette


{\vergleichskette


{ P
}


{ \in }{ V( \operatorname{Id} \,(T))
}


{  }{
}


{    }{
}


{   }{
}


}
{}{}{.}


}


  \zwischenueberschrift{Das Radikal}


\inputdefinition


{}


{


Ein
\definitionsverweis {Ideal}{}{}
${\mathfrak a}$ in einem
\definitionsverweis {kommutativen Ring}{}{}
$R$ heißt \definitionswort {Radikal}{}
\zusatzklammer {oder \definitionswort {Radikalideal}{}} {} {,}
wenn Folgendes gilt: Falls


\mavergleichskette


{\vergleichskette


{ f^n
}


{ \in }{ {\mathfrak a}
}


{  }{
}


{    }{
}


{   }{
}


}
{}{}{}
ist für ein


\mavergleichskette


{\vergleichskette


{ n
}


{ \in }{ \N
}


{  }{
}


{    }{
}


{   }{
}


}
{}{}{,}
so ist bereits


\mavergleichskette


{\vergleichskette


{ f
}


{ \in }{ {\mathfrak a}
}


{  }{
}


{    }{
}


{   }{
}


}
{}{}{.}


}


\inputdefinition


{ }


{


Es sei $R$ ein
\definitionsverweis {kommutativer Ring}{}{}
und


\mavergleichskette


{\vergleichskette


{ {\mathfrak a}
}


{ \subseteq }{ R
}


{  }{
}


{    }{
}


{   }{
}


}
{}{}{}
ein
\definitionsverweis {Ideal}{}{.}
Dann nennt man die Menge


\mathdisp {\left\{  f \in R   \mid   \text{es gibt ein } r \text{ mit } f^r \in {\mathfrak a}         \right\}} {  }


das \definitionswort {Radikal}{} zu ${\mathfrak a}$. Es wird mit


\mathl{\operatorname{rad}  \left(    {\mathfrak a}   \right)}{} bezeichnet.


}


Das Radikal zu einem Ideal ist selbst ein Radikal und insbesondere ein Ideal.


\inputfaktbeweis


{Theorie der Radikale (kommutative Algebra)/Radikal zu einem Ideal/ist ein Ideal/Fakt}


{Lemma}


{}


{


\faktsituation {}


\faktvoraussetzung {Es sei $R$ ein
\definitionsverweis {kommutativer Ring}{}{}
und


\mavergleichskette


{\vergleichskette


{ {\mathfrak a}
}


{ \subseteq }{ R
}


{  }{
}


{    }{
}


{   }{
}


}
{}{}{}
ein
\definitionsverweis {Ideal}{}{.}}


\faktfolgerung {Dann ist das
\definitionsverweis {Radikal}{}{}
zu ${\mathfrak a}$ ein
\definitionsverweis {Radikalideal}{}{.}}


\faktzusatz {}


\faktzusatz {}


}


{


Wir zeigen zunächst, dass ein Ideal vorliegt. $0$ gehört offenbar zum Radikal und mit


\mavergleichskette


{\vergleichskette


{ f
}


{ \in }{ \operatorname{rad}  \left(    {\mathfrak a}   \right)
}


{  }{
}


{    }{
}


{   }{
}


}
{}{}{,}
sagen wir


\mavergleichskette


{\vergleichskette


{ f^r
}


{ \in }{ {\mathfrak a}
}


{  }{
}


{    }{
}


{   }{
}


}
{}{}{,}
ist auch


\mavergleichskette


{\vergleichskette


{ (af)^r
}


{ = }{ a^rf^r
}


{ \in }{ {\mathfrak a}
}


{    }{
}


{   }{
}


}
{}{}{,}
also gehört $af$ zum Radikal. Zur Summeneigenschaft seien


\mavergleichskette


{\vergleichskette


{ f,g
}


{ \in }{ \operatorname{rad}  \left(    {\mathfrak a}   \right)
}


{  }{
}


{    }{
}


{   }{
}


}
{}{}{}
mit


\mavergleichskette


{\vergleichskette


{ f^r
}


{ \in }{ {\mathfrak a}
}


{  }{
}


{    }{
}


{   }{
}


}
{}{}{}
und


\mavergleichskette


{\vergleichskette


{ g^s
}


{ \in }{ {\mathfrak a}
}


{  }{
}


{    }{
}


{   }{
}


}
{}{}{.}
Dann ist


\mavergleichskettealign


{\vergleichskettealign


{ (f+g)^{r+s}
}


{ =} { \sum_{i+j = r+s} \binom{r+s}{i}  f^{i}g^{j}
}


{ =} { \sum_{i+j = r+s,\, i <r } \binom{r+s}{i}  f^{i}g^{j}  +\sum_{i+j = r+s,\, i \geq r} \binom{r+s}{i}  f^{i}g^{j}
}


{ \in} { {\mathfrak a}
}


{ } {
}


}
{}
{}{.}
Es sei nun


\mavergleichskette


{\vergleichskette


{ f^k
}


{ \in }{ \operatorname{rad}  \left(    {\mathfrak a}   \right)
}


{  }{
}


{    }{
}


{   }{
}


}
{}{}{.}
Dann ist


\mavergleichskette


{\vergleichskette


{ (f^k)^r
}


{ = }{ f^{kr}
}


{ \in }{ {\mathfrak a}
}


{    }{
}


{   }{
}


}
{}{}{,}
also


\mavergleichskette


{\vergleichskette


{ f
}


{ \in }{ \operatorname{rad}  \left(    {\mathfrak a}   \right)
}


{  }{
}


{    }{
}


{   }{
}


}
{}{}{.}


}


\inputfaktbeweis


{Affine Varietäten/Verschwindungsideal zu Teilmenge/ist Radikal/Fakt}


{Lemma}


{}


{


\faktsituation {}


\faktvoraussetzung {Es sei


\mavergleichskette


{\vergleichskette


{ T
}


{ \subseteq }{ { {\mathbb A}_{  K }^{  n   } }
}


{  }{
}


{    }{
}


{   }{
}


}
{}{}{}
eine Teilmenge.}


\faktfolgerung {Dann ist das
\definitionsverweis {Verschwindungsideal}{}{}
zu $T$ ein
\definitionsverweis {Radikal}{}{.}}


\faktzusatz {}


\faktzusatz {}


}


{


Es sei


\mavergleichskette


{\vergleichskette


{ F
}


{ \in }{ K[X_1 , \ldots ,  X_n]
}


{  }{
}


{    }{
}


{   }{
}


}
{}{}{}
ein Polynom, und sei


\mavergleichskette


{\vergleichskette


{ F^s
}


{ \in }{ \operatorname{Id} \,(T)
}


{  }{
}


{    }{
}


{   }{
}


}
{}{}{.}
Dann ist


\mavergleichskette


{\vergleichskette


{ F^s(P)
}


{ = }{ 0
}


{  }{
}


{    }{
}


{   }{
}


}
{}{}{}
für alle


\mavergleichskette


{\vergleichskette


{ P
}


{ \in }{ T
}


{  }{
}


{    }{
}


{   }{
}


}
{}{}{.}
Dann ist aber auch


\mavergleichskette


{\vergleichskette


{ F(P)
}


{ = }{ 0
}


{  }{
}


{    }{
}


{   }{
}


}
{}{}{}
für alle


\mavergleichskette


{\vergleichskette


{ P
}


{ \in }{ T
}


{  }{
}


{    }{
}


{   }{
}


}
{}{}{,}
also


\mavergleichskette


{\vergleichskette


{ F
}


{ \in }{ \operatorname{Id} \,(T)
}


{  }{
}


{    }{
}


{   }{
}


}
{}{}{.}


}


Wir werden später sehen, dass über einem algebraisch abgeschlossenen Körper sich Radikale und algebraische Nullstellengebilde entsprechen. Das ist der Inhalt des \stichwort {Hilbertschen Nullstellensatzes} {.}
